\documentclass[12pt]{article}

\usepackage{amsmath, amssymb, amsthm}
\usepackage{geometry}
\usepackage{accessibility}
\usepackage{listings}
\usepackage{xcolor}

\geometry{margin=1in}

\title{Rust Implementation of LoraHub's Mathematical Core}
\author{Frederick de Ruiter}
\date{\today}

\lstset{
    basicstyle=\ttfamily\small,
    keywordstyle=\color{blue},
    stringstyle=\color{red},
    commentstyle=\color{green!50!black},
    breaklines=true,
    frame=single
}

\begin{document}

\maketitle

\begin{abstract}
    This paper details the Rust implementation of the core mathematical principles behind LoraHub, a framework for the dynamic composition of Low-Rank Adaptation (LoRA) modules. We translate the mathematical concepts from the original LoraHub paper into a robust and efficient Rust implementation, leveraging the \texttt{nalgebra} library for linear algebra operations. The focus is on the two primary components: the weighted combination of LoRA tensors and the formulation of the objective function used for optimization.
\end{abstract}

\section{Introduction}

Low-Rank Adaptation (LoRA) has emerged as an efficient technique for fine-tuning large pre-trained models. LoraHub extends this by proposing a method to combine multiple specialized LoRA modules to achieve generalization on new, unseen tasks without additional training \cite{huang2023lorahub}. The central idea is to find an optimal linear combination of the weights of existing LoRA modules.

The core of LoraHub is a mathematical optimization problem. Given a set of LoRA modules (represented by their state dictionaries, which are collections of tensors), the goal is to find a set of scalar weights that, when used to combine the modules, minimizes an objective function on a few examples from a new task.

This paper documents a Rust implementation of the mathematical functions required for this process.

\section{Mathematical Formulation}

The LoraHub algorithm \cite{huang2023lorahub} can be broken down into two key mathematical components:

\begin{enumerate}
    \item \textbf{LoRA Combination:} The combined LoRA module is a weighted sum of the individual modules. If we let $S_i$ be the state dictionary (a set of tensors) for the $i$-th LoRA module and $w_i$ be its corresponding weight, the final combined state dictionary $S_f$ is given by:
    \begin{equation}
        S_f = \sum_{i=1}^{N} w_i S_i
    \end{equation}
    This operation is performed for each tensor within the state dictionaries.

    \item \textbf{Objective Function:} The optimal weights are found by minimizing an objective function. This function is typically the sum of a loss term and a regularization term:
    \begin{equation}
        \text{Objective}(w) = \text{Loss}(S_f) + \lambda R(w)
    \end{equation}
    where $\text{Loss}(S_f)$ is the model's loss on the new task using the combined LoRA, $R(w)$ is a regularization function on the weights, and $\lambda$ is a regularization parameter. LoraHub uses L1 regularization to encourage sparsity.
\end{enumerate}

\section{Rust Implementation}

We implement these concepts in Rust, creating a new module within the \texttt{math\_explorer} project. This implementation relies on the \texttt{nalgebra} library for its linear algebra capabilities \cite{nalgebra}.

\subsection{Combining LoRA Modules}

The \texttt{combine\_loras} function takes a slice of LoRA state dictionaries and a slice of weights and returns their weighted sum. A \texttt{LoraStateDict} is represented as a \texttt{HashMap<String, DMatrix<f64>>}.

\begin{lstlisting}
pub fn combine_loras(loras: &[LoraStateDict], weights: &[f64]) -> Result<LoraStateDict, &'static str> {
    if loras.is_empty() || weights.is_empty() {
        return Err("Input loras or weights cannot be empty.");
    }
    if loras.len() != weights.len() {
        return Err("The number of loras and weights must be the same.");
    }

    let first_lora = &loras[0];
    let mut final_state_dict: LoraStateDict = first_lora
        .iter()
        .map(|(key, tensor)| (key.clone(), tensor * weights[0]))
        .collect();

    for (i, lora_state_dict) in loras.iter().enumerate().skip(1) {
        let keys: Vec<String> = final_state_dict.keys().cloned().collect();
        for key in keys {
            if let Some(tensor) = lora_state_dict.get(&key) {
                if let Some(final_tensor) = final_state_dict.get_mut(&key) {
                    *final_tensor += tensor * weights[i];
                }
            }
        }
    }

    Ok(final_state_dict)
}
\end{lstlisting}

\subsection{Objective Function}

The objective function is implemented by \texttt{calculate\_objective\_score}, which combines a (mock) loss value with the L1 regularization term.

\begin{lstlisting}
pub fn l1_regularization(weights: &[f64], alpha: f64) -> f64 {
    if weights.is_empty() {
        return 0.0;
    }
    let sum_abs = weights.iter().map(|w| w.abs()).sum::<f64>();
    alpha * sum_abs / (weights.len() as f64)
}

pub fn calculate_objective_score(weights: &[f64], mock_loss: f64, alpha: f64) -> f64 {
    let regularization_term = l1_regularization(weights, alpha);
    mock_loss + regularization_term
}
\end{lstlisting}

\section{Conclusion}

We have successfully implemented the mathematical core of the LoraHub framework in Rust. This implementation provides a foundation for building a full LoraHub-style model composition system in Rust. The code is modular, tested, and integrated into the \texttt{math\_explorer} project, aligning with its goal of providing a library of mathematical algorithms. The next steps would involve integrating this mathematical core with a machine learning framework in Rust to replicate the full behavior of LoraHub.

\bibliographystyle{plain}
\bibliography{references}

\end{document}
